%%=============================================================================
%% Huidige situatie (AS-IS)
%%=============================================================================

\chapter{\IfLanguageName{dutch}{Huidige situatie}{Current situation}}%
\label{ch:huidige-situatie}

Dit hoofdstuk analyseert hoe CTF-events vandaag worden georganiseerd op eigen infrastructuur. Het brengt de problemen en beperkingen van de huidige aanpak in kaart als vertrekpunt voor de gewenste oplossing.

%%----------------------------------------------------------------------
\section{Organisatie van CTF-events op eigen infrastructuur}%
\label{sec:asis-organisatie}

% TODO: Beschrijf hoe CTF-events vandaag typisch worden georganiseerd op
% on-premise infrastructuur. Welke tools en processen worden gebruikt?
% Verwijs naar gesprekken met organisatoren of bestaande documentatie.
% Bronverwijzingen: \autocite{CTFd}, \autocite{Singier2025}

%%----------------------------------------------------------------------
\section{Manuele container deployments}%
\label{sec:asis-manueel}

% TODO: Bespreek hoe challenges vandaag worden uitgerold:
% - handmatige docker run / docker-compose commando's
% - scripts zonder foutafhandeling
% - geen geautomatiseerd opruimen van gestopte containers
% Welke problemen en beperkingen zijn hier aan verbonden?

%%----------------------------------------------------------------------
\section{Overzicht van bestaande tools en hun tekortkomingen}%
\label{sec:asis-tools}

% TODO: Geef een overzicht van de tools die vandaag gebruikt worden en
% beschrijf hun beperkingen:
% - CTFd zonder plugin: enkel statische challenges
% - CTFd-Whale: \autocite{CTFdWhale} - bestaande plugin, maar beperkingen?
% - Handmatige scripts: niet schaalbaar, geen isolatiegaranties
% - Hosted platforms (bijv. CTFd cloud): niet self-hosted, kostprijs

%%----------------------------------------------------------------------
\section{Conclusie: wat ontbreekt er?}%
\label{sec:asis-conclusie}

% TODO: Formuleer op basis van de analyse welke elementen ontbreken in de
% huidige aanpak. Dit vormt de directe input voor de TO-BE beschrijving.
% Denk aan: automatisering, isolatie, resource-beheer, schaalbaarheid,
% gebruiksgemak voor organisatoren.
